\def\handoutsay{false}\def\handoutfill{false}
\documentclass[12pt]{article}
\newcommand{\handoutrunning}{Lecture 2: Gauss Elimination, RREF, and Inverses}
\providecommand{\handoutsay}{true}
\usepackage[T1]{fontenc}
\usepackage{lmodern,microtype}
\usepackage[a4paper,margin=15mm,headheight=14pt,headsep=5mm]{geometry}
\usepackage{amsmath,amssymb,amsthm,mathtools}
\usepackage{enumitem,booktabs,needspace}
\usepackage{pifont}
\usepackage[dvipsnames]{xcolor}
\usepackage{fancyhdr}
\usepackage{ifthen}
\usepackage[hidelinks]{hyperref}
\hypersetup{pdfauthor={Chang-Ye Tu}}
\pagestyle{fancy}
\fancyhf{}
\fancyhead[L]{\small Linear Algebra}
\fancyhead[R]{\small\handoutrunning}
\fancyfoot[C]{\small\thepage}
\renewcommand{\headrulewidth}{0.3pt}
\setlength{\parindent}{0pt}
\setlength{\parskip}{4pt plus 1pt minus 1pt}
\setlength{\emergencystretch}{1.5em}
\usepackage{etoolbox}
\newcommand{\handoutdisplayskips}{\setlength{\abovedisplayskip}{7pt plus 2pt minus 2pt}%
  \setlength{\belowdisplayskip}{7pt plus 2pt minus 2pt}}
\apptocmd{\normalsize}{\handoutdisplayskips}{}{\PackageWarning{handout}{could not patch normalsize}}
\handoutdisplayskips
\setlist[enumerate]{label=(\alph*),leftmargin=*,itemsep=2pt,topsep=3pt}
\setlist[itemize]{leftmargin=*,itemsep=2pt,topsep=3pt}
\newcommand{\R}{\mathbb R}
\newcommand{\C}{\mathbb C}
\newcommand{\iu}{\mathrm{i}}
\newcommand{\contradiction}{\mathord{\text{\large\ding{56}}}}
\newcommand{\one}{\mathbf 1}
\DeclareMathOperator{\tr}{tr}
\DeclareMathOperator{\col}{col}
\DeclareMathOperator{\diag}{diag}
\DeclareMathOperator{\rank}{rank}
\DeclareMathOperator\prb{\mathsf{P}}
\DeclareMathOperator\expc{\mathsf{E}}
\DeclareMathOperator\var{var}
\DeclareMathOperator\cov{cov}
\DeclareMathOperator\cor{cor}
\DeclareMathOperator\argmin{argmin}
\DeclareMathOperator\argmax{argmax}
\DeclareMathOperator{\Span}{span}
\theoremstyle{definition}
\newtheorem{theorem}{Theorem}
\newtheorem{proposition}[theorem]{Proposition}
\newtheorem{lemma}[theorem]{Lemma}
\newtheorem{corollary}[theorem]{Corollary}
\newtheorem{definition}[theorem]{Definition}
\newtheorem{example}[theorem]{Example}
\newtheorem{remark}[theorem]{Remark}
\newtheorem{fact}{Fact}
\renewcommand{\thefact}{F\arabic{fact}}
\AddToHook{cmd/subsection/before}{\par\Needspace{9\baselineskip}}
\AddToHook{env/theorem/before}{\par\Needspace{4\baselineskip}}
\AddToHook{env/proposition/before}{\par\Needspace{4\baselineskip}}
\AddToHook{env/lemma/before}{\par\Needspace{4\baselineskip}}
\AddToHook{env/definition/before}{\par\Needspace{4\baselineskip}}
\AddToHook{env/example/before}{\par\Needspace{4\baselineskip}}
\definecolor{saybg}{gray}{0.86}
\NewDocumentEnvironment{say}{+b}
  {\ifthenelse{\equal{\handoutsay}{true}}{\par\Needspace{2\baselineskip}\noindent\colorbox{saybg}{\parbox{\dimexpr\linewidth-2\fboxsep\relax}%
   {\small\sffamily\ding{43}\enspace\textbf{Say.}\enspace#1}}\par}{}}
  {}
\providecommand{\handoutfill}{false}
\newsavebox{\boardbox}
\newlength{\boardht}
\definecolor{boardrule}{gray}{0.55}
\makeatletter
\newcommand{\boardstart}{\par\if@inlabel\leavevmode\par\fi}
\makeatother
\NewDocumentEnvironment{board}{o +b}
  {\boardstart\IfValueT{#1}{\noindent#1\par\nopagebreak}%
   \noindent\begin{lrbox}{\boardbox}\begin{minipage}{\linewidth}#2\end{minipage}\end{lrbox}%
   \setlength{\boardht}{\dimexpr\ht\boardbox+\dp\boardbox\relax}%
   \ifthenelse{\equal{\handoutfill}{true}}%
     {\setlength{\boardht}{1.4\boardht}%
      {\setlength{\fboxsep}{0pt}\setlength{\fboxrule}{0.4pt}\color{boardrule}%
       \framebox[\dimexpr\linewidth-2\fboxrule\relax]{\rule{0pt}{\boardht}}}}%
     {\ifthenelse{\equal{\handoutsay}{true}}%
        {\llap{\color{boardrule}\rule[-\dp\boardbox]{1.5pt}{\boardht}\hspace{4pt}}}{}%
      \usebox{\boardbox}}%
   \par}
  {}
\newcommand{\blank}[1]{\ifthenelse{\equal{\handoutfill}{true}}{\underline{\phantom{#1}}}{#1}}
\newcommand{\lec}[1]{Lecture~#1}
\newcommand{\handouttitle}[2]{%
  \begin{center}
  {\LARGE\sffamily\bfseries #2\par}\vspace{3pt}
  {\sffamily Linear Algebra \textperiodcentered\ \lec{#1}\par}
  \ifthenelse{\equal{\handoutsay}{true}}{\vspace{3pt}{\small\sffamily\colorbox{saybg}{\ding{43}\enspace\textbf{Say.}}\enspace
    teacher's copy: the shaded boxes are not in the students' copy.\par}}{}
  \ifthenelse{\equal{\handoutfill}{true}}{\vspace{3pt}{\small\sffamily Fill-in copy: complete the framed spaces in
    class.\par}}{}
  \end{center}}

\hypersetup{pdftitle={Linear Algebra, Lecture 2: Gauss Elimination, RREF, and Inverses}}
\newcommand{\boardtop}{\vspace*{-\abovedisplayskip}}
\DeclareMathOperator{\row}{row}
\begin{document}
\handouttitle{2}{Gauss Elimination, RREF, and Inverses}

\textbf{Last time} (\lec{1}): products of matrices and their rules; the columns and rows of a product; the transpose.
\(e_i\) is column \(i\) of the identity matrix (\(I_m\) or \(I_n\), as the product requires).

\subsection*{1. Linear systems and row operations}

\begin{definition}\label{b2.1}
Let \(A=(a_{ij})\in\R^{m\times n}\) and \(b\in\R^m\). The system of \(m\) linear equations
\(a_{i1}x_1+\dots+a_{in}x_n=b_i\) (\(i=1,\dots,m\)) in \(n\) unknowns is \(Ax=b\) (\lec{1},
Definition~2). \(A\) is the \textbf{coefficient matrix}, \([A\mid b]\in\R^{m\times(n+1)}\) the
\textbf{augmented matrix}. A \textbf{solution} is an \(x\in\R^n\) with \(Ax=b\). The system is \textbf{consistent} if
it has a solution, \textbf{inconsistent} otherwise; \textbf{homogeneous} if \(b=0\), and then \(x=0\) is a solution,
the \textbf{trivial solution}.
\end{definition}

\begin{definition}\label{b2.2}
The \textbf{elementary row operations} on a matrix with \(m\) rows are, for rows \(r\ne s\):
(1) \(R_r\leftarrow cR_r\) (\(c\ne0\)): multiply row \(r\) by \(c\); (2) \(R_r\leftarrow R_r+cR_s\) (\(c\in\R\)): add
\(c\) times row \(s\) to row \(r\); (3) \(R_r\leftrightarrow R_s\): interchange rows \(r\) and \(s\). Several operations
on one arrow are performed in the order written (upper label first). \(B\) is \textbf{row-equivalent} to \(A\) if
finitely many elementary row operations take \(A\) to \(B\).
\end{definition}
An operation computes each column of the new matrix from the same column of the old one: operations that take
\([A\mid b]\) to \([B\mid z]\) take \(A\) to \(B\), and operations that take \(A\) to \(R\) take \([A\mid0]\) to
\([R\mid0]\).

\Needspace{11\baselineskip}
\begin{theorem}\label{b2.3}
\begin{enumerate}
\item Each elementary row operation is undone by one of the same type: \(R_r\leftarrow cR_r\) by
\(R_r\leftarrow c^{-1}R_r\), \(R_r\leftarrow R_r+cR_s\) by \(R_r\leftarrow R_r-cR_s\), \(R_r\leftrightarrow R_s\) by
itself.
\item If row operations take \([A\mid b]\) to \([B\mid z]\), then \(Ax=b\) and \(Bx=z\) have the same solutions.
\end{enumerate}
\end{theorem}
\begin{proof}
\begin{board}[(a) Row by row; the other rows do not change:]
\boardtop
\begin{align*}
&c^{-1}(c\row_r)=\row_r\\
&(\row_r+c\row_s)-c\row_s=\row_r&&\text{(row \(s\) is unchanged)}\\
&\text{two interchanges of rows \(r\), \(s\) restore their order}
\end{align*}
\end{board}
\begin{board}[(b) It suffices to take one operation:]
\boardtop
\begin{align*}
Ax=b&\implies Bx=z&&\text{(each equation of \(Bx=z\) is one of \(Ax=b\), \(c\) times one,}\\
&&&\text{\phantom{(}or one plus \(c\) times another)}\\
Bx=z&\implies Ax=b&&\text{((a): an operation takes \([B\mid z]\) back to \([A\mid b]\))}
\end{align*}
\end{board}
\nopagebreak
\end{proof}

\subsection*{2. Reduced row echelon form}

\begin{definition}\label{b2.4}
A matrix \(R\) is in \textbf{reduced row echelon form} (\textbf{RREF}) if
\begin{enumerate}
\item the first nonzero entry of each nonzero row is \(1\), its \textbf{pivot};
\item each column that contains a pivot has all its other entries \(0\);
\item every zero row lies below every nonzero row;
\item if rows \(1,\dots,r\) are the nonzero rows and the pivot of row \(i\) is in column \(k_i\), then
\(k_1<\dots<k_r\).
\end{enumerate}
Columns \(k_1,\dots,k_r\) are the \textbf{pivot columns}; column \(k_i\) of \(R\) is \(e_i\) (by (a), (b)). A matrix
that satisfies (a), (c), (d) is in \textbf{echelon form}; then every entry below a pivot is \(0\) (by (c), (d)).
\end{definition}
For example, \(I_n\) and \(0\) are in RREF, and so is
\(R=\left(\begin{smallmatrix}0&1&-3&0&2\\0&0&0&1&2\\0&0&0&0&0\end{smallmatrix}\right)\) (\(r=2\), \(k_1=2\),
\(k_2=4\)). \(\left(\begin{smallmatrix}1&0&0&0\\0&1&-1&0\\0&0&1&0\end{smallmatrix}\right)\) is in echelon form, not in
RREF (column \(3\) violates (b)); \(\left(\begin{smallmatrix}0&2&1\\1&0&-3\\0&0&0\end{smallmatrix}\right)\) violates (a)
and (d).

\begin{definition}[Gauss--Jordan elimination]\label{b2.5}
Let \(M\) be any matrix. Start with \(i=1\).
\begin{enumerate}[label=\textit{Step \arabic*.},leftmargin=*,itemsep=0pt]
\item Find the first column with a nonzero entry in rows \(i,i+1,\dots\) (the next pivot column); if there is none,
go to Step 5.
\item If row \(i\) has \(0\) in this column, interchange it with a lower row that has a nonzero entry there.
\item Multiply row \(i\) by the reciprocal of this entry, so that the pivot is \(1\).
\item Add multiples of row \(i\) to the rows below it to make every entry below the pivot \(0\). Increase \(i\) by
\(1\); return to Step 1.
\item (\textbf{Backward elimination}.) For each pivot, from the last to the first, add multiples of its row to the
rows above it to make every entry above the pivot \(0\).
\end{enumerate}
Steps 1--4 (\textbf{forward elimination}) produce an echelon form, Step 5 a matrix in RREF:
\begin{itemize}
\item When Step 1 finds the pivot column of row \(i\), rows \(i,i+1,\dots\) are \(0\) in every column before it, and
Step 4 makes rows \(i+1,i+2,\dots\) \(0\) in it: each nonzero row begins with its pivot \(1\), the pivot columns are
taken from left to right, the rows left when Step 1 finds none are zero; an echelon form, pivots in columns
\(k_1<\dots<k_r\).
\item In Step 5, when the pivot of row \(i\) is processed, row \(i\) is \(0\) before column \(k_i\) and in columns
\(k_{i+1},\dots,k_r\) (done): the multiples of row \(i\) change no entry before column \(k_i\) and no other pivot
column. In the end every column \(k_i\) is \(e_i\).
\item Step 3 may be postponed, to the end of forward elimination or to the very end (the pivots become \(1\) only
then); this often avoids fractions.
\end{itemize}
So every matrix \(M\) has an \textbf{RREF}: a matrix in RREF that is row-equivalent to \(M\).
\end{definition}

\Needspace{11.5\baselineskip}
\begin{example}\label{b2.6}
Solve \(x_1+x_2+x_3-2x_4=0\), \(x_1+x_2+2x_3-3x_4=0\), \(2x_1+x_2+5x_3-4x_4=0\).
\end{example}
\textit{Solution.}
\begin{board}[Row-reduce the coefficient matrix \(A\) (the column \(b=0\) never changes; not written):]
\[
\begin{gathered}
\begin{pmatrix}1&1&1&-2\\1&1&2&-3\\2&1&5&-4\end{pmatrix}
\xrightarrow[R_3\leftarrow R_3-2R_1]{R_2\leftarrow R_2-R_1}
\begin{pmatrix}1&1&1&-2\\0&0&1&-1\\0&-1&3&0\end{pmatrix}
\xrightarrow{R_2\leftrightarrow R_3}
\begin{pmatrix}1&1&1&-2\\0&-1&3&0\\0&0&1&-1\end{pmatrix}
\\
\xrightarrow{R_2\leftarrow-R_2}
\begin{pmatrix}1&1&1&-2\\0&1&-3&0\\0&0&1&-1\end{pmatrix}
\xrightarrow[R_1\leftarrow R_1-R_3]{R_2\leftarrow R_2+3R_3}
\begin{pmatrix}1&1&0&-1\\0&1&0&-3\\0&0&1&-1\end{pmatrix}
\xrightarrow{R_1\leftarrow R_1-R_2}
\begin{pmatrix}1&0&0&2\\0&1&0&-3\\0&0&1&-1\end{pmatrix}
\end{gathered}
\]
\end{board}
Entry \((2,2)\) is \(0\) after the first arrow: Step 2 interchanges rows \(2\), \(3\); Step 3 multiplies row \(2\) by
\(-1\). The first matrix of the second line is an echelon form.
\begin{board}[The last matrix \(R\) is in RREF, pivot columns \(1,2,3\); \(Rx=0\) has the solutions of \(Ax=0\)
(Theorem~\ref{b2.3}(b)):]
\boardtop
\begin{align*}
&x_1+2x_4=0,\quad x_2-3x_4=0,\quad x_3-x_4=0;\qquad x_4=t\text{ free}:\quad x=t(-2,3,1,1)'\quad(t\in\R)
\end{align*}
\end{board}
Check: \(-2+3+1-2=0\), \(-2+3+2-3=0\), \(-4+3+5-4=0\).

\Needspace{13.5\baselineskip}
\begin{proposition}\label{b2.7}
Let row operations take \([A\mid b]\) (\(A\in\R^{m\times n}\)) to \([R\mid z]\), \(R\) in RREF with nonzero rows
\(1,\dots,r\) and pivot columns \(k_1<\dots<k_r\); \(x_{k_1},\dots,x_{k_r}\) are the \textbf{pivot unknowns}, the
other \(n-r\) unknowns the \textbf{free unknowns}.
\begin{enumerate}
\item \(Ax=b\) is consistent \(\iff\) \(z_i=0\) for every \(i>r\) (no row \((0\ \cdots\ 0\mid z_i)\), \(z_i\ne0\)).
\item If so, all solutions: the free unknowns arbitrary, \(x_{k_i}=z_i-\sum_{j\text{ free}}R_{ij}x_j\)
(\(1\leqslant i\leqslant r\)).
\end{enumerate}
\end{proposition}
\begin{proof}
\(Ax=b\) and \(Rx=z\) have the same solutions (Theorem~\ref{b2.3}(b)); column \(k_l\) of \(R\) is \(e_l\)
(Definition~\ref{b2.4}).
\begin{board}[(a) Equation \(i>r\) of \(Rx=z\) reads \(0=z_i\):]
\boardtop
\begin{align*}
&z_i\ne0\text{ for some }i>r\implies\text{no solution}\\
&z_i=0\text{ for all }i>r:\ \text{take the free unknowns }0\text{ and }x_{k_l}=z_l\ (l\leqslant r)\\
&\implies Rx=\textstyle\sum_{l=1}^rz_l\col_{k_l}(R)=\sum_{l=1}^rz_le_l=z&&\text{(\lec{1}, Theorem~4(b))}
\end{align*}
\end{board}
\begin{board}[(b) \(x_{k_i}\) occurs only in equation \(i\), with coefficient \(1\); equations \(i>r\) read \(0=0\):]
\boardtop
\begin{align*}
&\text{equation }i\leqslant r:\quad x_{k_i}+\textstyle\sum_{j\text{ free}}R_{ij}x_j=z_i\\
&\implies\text{each choice of the free unknowns gives exactly one solution, and every solution arises so}
\end{align*}
\end{board}
\nopagebreak
\end{proof}

\Needspace{15\baselineskip}
\begin{example}\label{b2.10}
Let \(A=\left(\begin{smallmatrix}1&2&0&1\\2&3&1&1\\1&1&1&1\\2&2&2&1\end{smallmatrix}\right)\). For which
\(y=(y_1,y_2,y_3,y_4)'\) is \(Ax=y\) consistent? Solve it for \(y=(4,7,4,7)'\) and for \(y=(1,0,0,0)'\).
\end{example}
\textit{Solution.}
\begin{board}[Forward elimination of \([A\mid y]\), keeping \(y_1,\dots,y_4\) as letters:]
\[
\begin{gathered}
\left(\begin{array}{cccc|c}1&2&0&1&y_1\\2&3&1&1&y_2\\1&1&1&1&y_3\\2&2&2&1&y_4\end{array}\right)
\xrightarrow[R_4\leftarrow R_4-2R_1]{R_2\leftarrow R_2-2R_1,\ R_3\leftarrow R_3-R_1}
\left(\begin{array}{cccc|c}1&2&0&1&y_1\\0&-1&1&-1&y_2-2y_1\\0&-1&1&0&y_3-y_1\\0&-2&2&-1&y_4-2y_1\end{array}\right)
\\
\xrightarrow{R_2\leftarrow-R_2}
\left(\begin{array}{cccc|c}1&2&0&1&y_1\\0&1&-1&1&2y_1-y_2\\0&-1&1&0&y_3-y_1\\0&-2&2&-1&y_4-2y_1\end{array}\right)
\xrightarrow[R_4\leftarrow R_4+2R_2]{R_3\leftarrow R_3+R_2}
\left(\begin{array}{cccc|c}1&2&0&1&y_1\\0&1&-1&1&2y_1-y_2\\0&0&0&1&y_1-y_2+y_3\\0&0&0&1&2y_1-2y_2+y_4\end{array}\right)
\end{gathered}
\]
\end{board}
\begin{board}[The last step of forward elimination (the left part is now an echelon form); then backward
elimination:]
\[
\begin{gathered}
\xrightarrow{R_4\leftarrow R_4-R_3}
\left(\begin{array}{cccc|c}1&2&0&1&y_1\\0&1&-1&1&2y_1-y_2\\0&0&0&1&y_1-y_2+y_3\\0&0&0&0&y_1-y_2-y_3+y_4\end{array}\right)
\xrightarrow[R_1\leftarrow R_1-R_3]{R_2\leftarrow R_2-R_3}
\left(\begin{array}{cccc|c}1&2&0&0&y_2-y_3\\0&1&-1&0&y_1-y_3\\0&0&0&1&y_1-y_2+y_3\\0&0&0&0&y_1-y_2-y_3+y_4\end{array}\right)
\\
\xrightarrow{R_1\leftarrow R_1-2R_2}
\left(\begin{array}{cccc|c}1&0&2&0&-2y_1+y_2+y_3\\0&1&-1&0&y_1-y_3\\0&0&0&1&y_1-y_2+y_3\\0&0&0&0&y_1-y_2-y_3+y_4\end{array}\right)
\end{gathered}
\]
\end{board}
\begin{board}[Read off (Proposition~\ref{b2.7}): the left part is an RREF of \(A\) (Definition~\ref{b2.2}),
\(r=3\), pivot columns \(1,2,4\), \(z_4=y_1-y_2-y_3+y_4\):]
\boardtop
\begin{align*}
&\text{consistent}\iff y_1-y_2-y_3+y_4=0;\ \text{then, with }x_3=c\text{ free }(c\in\R)\text{:}\\
&x_1=-2y_1+y_2+y_3-2c,\quad x_2=y_1-y_3+c,\quad x_3=c,\quad x_4=y_1-y_2+y_3\\
&y=(4,7,4,7)':\ 4-7-4+7=0,\quad x=(3,0,0,1)'+c(-2,1,1,0)'\\
&y=(1,0,0,0)':\ 1-0-0+0=1\ne0,\quad\text{no solution}
\end{align*}
\end{board}
Check: \(A(3,0,0,1)'=(4,7,4,7)'\); \(A(-2,1,1,0)'=0\), that is, \(\col_3(A)=2\col_1(A)-\col_2(A)\).

\Needspace{9\baselineskip}
\begin{theorem}\label{b2.11}
\begin{enumerate}
\item If \(A\in\R^{m\times n}\) with \(m<n\), then \(Ax=0\) has a nontrivial solution.
\item Let \(A\in\R^{n\times n}\) and \(R\) an RREF of \(A\). Then \(Ax=0\) has only the trivial solution \(\iff\)
\(R=I_n\).
\end{enumerate}
\end{theorem}
\begin{proof}
Let \(R\) be an RREF of \(A\) (Definition~\ref{b2.5}), with \(r\) nonzero rows; the operations take \([A\mid0]\) to
\([R\mid0]\) (Definition~\ref{b2.2}).
\begin{board}[(a) Proposition~\ref{b2.7}(b):]
\boardtop
\begin{align*}
r\leqslant m<n&\implies\text{a free unknown exists; it }=1,\text{ the others }=0\implies\text{a solution }x\ne0
\end{align*}
\end{board}
\begin{board}[(b) Theorem~\ref{b2.3}(b), Proposition~\ref{b2.7}(b); \lec{1}, Theorem~3(a):]
\boardtop
\begin{align*}
R=I&\implies Rx=x\implies x=0\text{ is the only solution}\\
\text{only }x=0&\implies\text{no free unknown (one would give \(x\ne0\))}\implies r=n\\
&\implies1\leqslant k_1<\dots<k_n\leqslant n\implies k_i=i,\ \col_i(R)=e_i\implies R=I&&\text{(Definition~\ref{b2.4})}
\end{align*}
\end{board}
\nopagebreak
\end{proof}

\subsection*{3. Inverses and elementary matrices}

\begin{definition}\label{b2.12}
\(A\in\R^{n\times n}\) is \textbf{invertible} if \(AB=BA=I_n\) for some \(B\in\R^{n\times n}\). Such \(B\) is unique:
\(BA=I\), \(AB_1=I\) \(\Rightarrow\) \(B=B(AB_1)=(BA)B_1=B_1\) (\lec{1}, Theorem~3(a)). It is the
\textbf{inverse} \(A^{-1}\) of \(A\). For example, if all \(d_i\ne0\), then
\(\diag(d_1,\dots,d_n)^{-1}=\diag(d_1^{-1},\dots,d_n^{-1})\): both products are \(\diag(1,\dots,1)=I\) (\lec{1},
Proposition~13(b)).
\end{definition}

\Needspace{11.5\baselineskip}
\begin{theorem}\label{b2.13}
Let \(A,B\in\R^{n\times n}\) be invertible.
\begin{enumerate}
\item \(A^{-1}\) is invertible and \((A^{-1})^{-1}=A\).
\item \(AB\) is invertible and \((AB)^{-1}=B^{-1}A^{-1}\); more generally \((A_1\cdots A_k)^{-1}=A_k^{-1}\cdots A_1^{-1}\)
for invertible \(A_1,\dots,A_k\in\R^{n\times n}\).
\item \(A'\) is invertible and \((A')^{-1}=(A^{-1})'\).
\end{enumerate}
\end{theorem}
\begin{proof}
\begin{board}[(a) Definition~\ref{b2.12}:]
\[
AA^{-1}=A^{-1}A=I\implies A\text{ is the inverse of }A^{-1}
\]
\end{board}
\begin{board}[(b) \lec{1}, Theorem~3(a):]
\boardtop
\begin{align*}
(AB)(B^{-1}A^{-1})&=A(BB^{-1})A^{-1}=AA^{-1}=I,\qquad(B^{-1}A^{-1})(AB)=B^{-1}(A^{-1}A)B=B^{-1}B=I\\
(A_1\cdots A_k)^{-1}&=\bigl((A_1\cdots A_{k-1})A_k\bigr)^{-1}=A_k^{-1}(A_1\cdots A_{k-1})^{-1}=\dots=A_k^{-1}\cdots A_1^{-1}
\end{align*}
\end{board}
\begin{board}[(c) \((AB)'=B'A'\), \(I'=I\) (\lec{1}, Theorem~8):]
\[
A'(A^{-1})'=(A^{-1}A)'=I'=I,\qquad(A^{-1})'A'=(AA^{-1})'=I
\]
\end{board}
\nopagebreak
\end{proof}

\Needspace{6\baselineskip}
\begin{corollary}\label{b2.23}
If \(A\in\R^{n\times n}\) is symmetric and invertible, then \(A^{-1}\) is symmetric.
\end{corollary}
\begin{proof}
\begin{board}[\(A'=A\) (\lec{1}, Definition~7); Theorem~\ref{b2.13}(c):]
\[
(A^{-1})'=(A')^{-1}=A^{-1}
\]
\end{board}
\nopagebreak
\end{proof}

\begin{definition}\label{b2.14}
An \textbf{elementary matrix} is \(E=\omega(I_m)\), obtained from \(I_m\) by one elementary row operation \(\omega\).
For \(r\ne s\): \(D_r(c)\) (\(c\ne0\)) is \(I_m\) with entry \((r,r)\) replaced by \(c\) (from \(R_r\leftarrow cR_r\));
\(T_{rs}(c)\) is \(I_m\) with entry \((r,s)\) replaced by \(c\) (from \(R_r\leftarrow R_r+cR_s\)); \(P_{rs}\) is \(I_m\)
with rows \(r\), \(s\) interchanged (from \(R_r\leftrightarrow R_s\)). For example, for \(m=4\):
\[
P_{23}=\begin{pmatrix}1&0&0&0\\0&0&1&0\\0&1&0&0\\0&0&0&1\end{pmatrix},\qquad
T_{31}(c)=\begin{pmatrix}1&0&0&0\\0&1&0&0\\c&0&1&0\\0&0&0&1\end{pmatrix},\qquad
D_3(c)=\begin{pmatrix}1&0&0&0\\0&1&0&0\\0&0&c&0\\0&0&0&1\end{pmatrix}
\]
\end{definition}

\Needspace{16\baselineskip}
\begin{theorem}\label{b2.15}
Let \(\omega\) be an elementary row operation on matrices with \(m\) rows and \(E=\omega(I_m)\).
\begin{enumerate}
\item \(\omega(A)=EA\) for every \(A\in\R^{m\times n}\).
\item \(E\) is invertible, and \(E^{-1}\) is the elementary matrix of the inverse operation:
\(D_r(c)^{-1}=D_r(c^{-1})\), \(T_{rs}(c)^{-1}=T_{rs}(-c)\), \(P_{rs}^{-1}=P_{rs}\).
\item If row operations \(\omega_1,\dots,\omega_k\) (in this order) take \(A\) to \(B\), then \(B=PA\) with
\(P=E_k\cdots E_2E_1\), \(E_j=\omega_j(I_m)\), and \(P\) is invertible.
\end{enumerate}
\end{theorem}
\begin{proof}
\begin{board}[(a) \(\row_i(EA)=\row_i(E)A\) (\lec{1}, Theorem~4(c)); an unchanged row of \(E\) is
\(e_i'\), and \(e_i'A=\row_i(A)\) (\lec{1}, Proposition~11(a)); \lec{1}, Theorem~3(a):]
\boardtop
\begin{align*}
R_r\leftarrow cR_r:&\quad\row_r(EA)=(ce_r')A=c\row_r(A)\\
R_r\leftarrow R_r+cR_s:&\quad\row_r(EA)=(e_r'+ce_s')A=\row_r(A)+c\row_s(A)\\
R_r\leftrightarrow R_s:&\quad\row_r(EA)=e_s'A=\row_s(A),\quad\row_s(EA)=e_r'A=\row_r(A)
\end{align*}
\end{board}
\begin{board}[(b) \(\tilde\omega\): the inverse operation (Theorem~\ref{b2.3}(a)), \(\tilde E=\tilde\omega(I)\); by (a):]
\[
E\tilde E=\omega(\tilde\omega(I))=I,\qquad\tilde EE=\tilde\omega(\omega(I))=I
\]
\end{board}
\begin{board}[(c) By (a) and \lec{1}, Theorem~3(a); Theorem~\ref{b2.13}(b) with (b):]
\[
B=E_k(\cdots(E_2(E_1A)))=(E_k\cdots E_1)A=PA,\qquad P\text{ invertible}
\]
\end{board}
\nopagebreak
\end{proof}

\subsection*{4. Deciding invertibility and computing inverses}

\begin{theorem}\label{b2.16}
Let \(A\in\R^{n\times n}\) and \(R\) an RREF of \(A\). The following are equivalent:
(i) \(A\) is invertible; (ii) \(R=I_n\); (iii) \(Ax=0\) has only the trivial solution; (iv) \(Ax=b\) has a solution
for every \(b\in\R^n\); (v) \(A\) is a product of elementary matrices. Moreover, if \(BA=I\) or \(AB=I\) for some
\(B\in\R^{n\times n}\), then \(A\) is invertible and \(B=A^{-1}\).
\end{theorem}
\begin{proof}
\(R=PA\), \(P=E_k\cdots E_1\) invertible (Theorem~\ref{b2.15}(c)).
\begin{itemize}
\item \(R=I\): \(A=P^{-1}(PA)=P^{-1}R=P^{-1}\) (\lec{1}, Theorem~3(a)) is \(E_1^{-1}\cdots E_k^{-1}\), a product of elementary matrices, hence invertible
(Theorems~\ref{b2.15}(b), \ref{b2.13}(b)); \(Ax=0\) has only the trivial solution (Theorem~\ref{b2.11}(b)); and
\(x=A^{-1}b\) solves \(Ax=b\).
\item \(R\ne I\): then \(R\) has a zero row, hence a zero last row (Definition~\ref{b2.4}(c)); otherwise \(r=n\),
\(1\leqslant k_1<\dots<k_n\leqslant n\) forces \(k_i=i\), \(\col_i(R)=e_i\) for every \(i\) (Definition~\ref{b2.4}),
and \(R=I\).
\item \(R\ne I\): \(Ax=0\) for some \(x\ne0\) (Theorem~\ref{b2.11}(b)), so \(A\) has no inverse (else
\(x=A^{-1}Ax=0\) \(\contradiction\)), and \(A\) is not a product of elementary matrices (such a product is invertible: Theorems~\ref{b2.15}(b),
\ref{b2.13}(b)).
For \(b=P^{-1}e_n\): \(Ax=b\iff Rx=PAx=Pb=e_n\), whose last equation reads \(0=1\); so (iv) fails.
\item So each of (i)--(v) holds exactly when \(R=I\).
\end{itemize}
\begin{board}[Last sentence; (iii) \(\Rightarrow\) (i):]
\boardtop
\begin{align*}
BA=I&:\ Ax=0\implies x=BAx=0\implies A\text{ invertible},\ B=BAA^{-1}=A^{-1}\\
AB=I&:\ B\text{ invertible},\ B^{-1}=A\ \text{(case just shown)}\implies A^{-1}=(B^{-1})^{-1}=B
&&\text{(Theorem~\ref{b2.13}(a))}
\end{align*}
\end{board}
\nopagebreak
\end{proof}

\Needspace{10.5\baselineskip}
\begin{corollary}[Computing \(A^{-1}\)]\label{b2.17}
Let \(A\in\R^{n\times n}\); row-reduce \([A\mid I_n]\) by Gauss--Jordan elimination.
\begin{enumerate}
\item If the left half becomes \(I_n\), the right half is \(A^{-1}\).
\item If a zero row appears in the left half, \(A\) is not invertible.
\end{enumerate}
\end{corollary}
\begin{proof}
\begin{board}[(a) \(E\): the product of the elementary matrices used (Theorem~\ref{b2.15}(c)); \(\col_j(EM)=E\col_j(M)\)
(\lec{1}, Theorem~4(a)):]
\[
[I\mid B]=E[A\mid I]=[EA\mid E]\implies EA=I\implies B=E=A^{-1}\qquad\text{(Theorem~\ref{b2.16}, last sentence)}
\]
\end{board}
\begin{board}[(b) The left half is \(EA\), \(E\) invertible (Theorem~\ref{b2.15}(c)); \(\row_i(EAX)=\row_i(EA)X\)
(\lec{1}, Theorem~4(c)):]
\boardtop
\begin{align*}
&\row_i(EA)=0\implies\row_i(EAX)=0\implies EAX\ne I_n\ \text{for all }X\in\R^{n\times n}\\
&\implies EA\text{ not invertible}\\
&A\text{ invertible}\implies EA\text{ invertible}\ \contradiction&&\text{(Theorem~\ref{b2.13}(b))}
\end{align*}
\end{board}
\nopagebreak
\end{proof}

\Needspace{14.5\baselineskip}
\begin{proposition}\label{b2.18}
Let \(A=\begin{pmatrix}a&b\\c&d\end{pmatrix}\in\R^{2\times2}\). Then \(A\) is invertible \(\iff\) \(ad-bc\ne0\), and
then
\[
A^{-1}=\frac1{ad-bc}\begin{pmatrix}d&-b\\-c&a\end{pmatrix}
\]
\end{proposition}
\begin{proof}
Row-reduce \([A\mid I_2]\) and read the result by Corollary~\ref{b2.17}; \(\Delta=ad-bc\).
\begin{board}[\(a\ne0\) (\(d-\frac{cb}a=\frac\Delta a\); \(\frac1a+\frac{bc}{a\Delta}=\frac{\Delta+bc}{a\Delta}
=\frac d\Delta\)):]
\[
\begin{gathered}
\left(\begin{array}{cc|cc}a&b&1&0\\c&d&0&1\end{array}\right)
\xrightarrow{R_2\leftarrow R_2-\frac caR_1}
\left(\begin{array}{cc|cc}a&b&1&0\\0&\frac\Delta a&-\frac ca&1\end{array}\right)
\qquad\Delta=0\implies\text{zero row}\\
\Delta\ne0:\
\xrightarrow[R_2\leftarrow\frac a\Delta R_2]{R_1\leftarrow\frac1aR_1}
\left(\begin{array}{cc|cc}1&\frac ba&\frac1a&0\\0&1&-\frac c\Delta&\frac a\Delta\end{array}\right)
\xrightarrow{R_1\leftarrow R_1-\frac baR_2}
\left(\begin{array}{cc|cc}1&0&\frac d\Delta&-\frac b\Delta\\0&1&-\frac c\Delta&\frac a\Delta\end{array}\right)
\end{gathered}
\]
\(\Delta=0\): not invertible (Corollary~\ref{b2.17}(b)); \(\Delta\ne0\): the right half is the stated \(A^{-1}\).
\end{board}
\begin{board}[\(a=0\), so \(\Delta=-bc\); \(c=0\): \(Ae_1=\col_1(A)=0\) (\lec{1}, Proposition~11(a)), not invertible
(Theorem~\ref{b2.16}, (i)\(\Leftrightarrow\)(iii)); \(c\ne0\):]
\[
\begin{gathered}
\left(\begin{array}{cc|cc}0&b&1&0\\c&d&0&1\end{array}\right)
\xrightarrow{R_1\leftrightarrow R_2}
\left(\begin{array}{cc|cc}c&d&0&1\\0&b&1&0\end{array}\right)
\quad b=0\implies\text{zero row, not invertible (Corollary~\ref{b2.17}(b))}\\
b\ne0:\
\xrightarrow[R_2\leftarrow\frac1bR_2]{R_1\leftarrow\frac1cR_1}
\left(\begin{array}{cc|cc}1&\frac dc&0&\frac1c\\0&1&\frac1b&0\end{array}\right)
\xrightarrow{R_1\leftarrow R_1-\frac dcR_2}
\left(\begin{array}{cc|cc}1&0&-\frac d{bc}&\frac1c\\0&1&\frac1b&0\end{array}\right)
\end{gathered}
\]
The right half is \(\frac1\Delta\left(\begin{smallmatrix}d&-b\\-c&0\end{smallmatrix}\right)\), the stated \(A^{-1}\)
for \(a=0\).
\end{board}
So \(A\) is invertible exactly when \(\Delta\ne0\) (for \(a=0\): \(b\ne0\), \(c\ne0\)), with the stated inverse
(Corollary~\ref{b2.17}(a)).
\end{proof}
Check by multiplication: \(\left(\begin{smallmatrix}a&b\\c&d\end{smallmatrix}\right)
\left(\begin{smallmatrix}d&-b\\-c&a\end{smallmatrix}\right)=\left(\begin{smallmatrix}ad-bc&0\\0&ad-bc\end{smallmatrix}\right)
=\left(\begin{smallmatrix}d&-b\\-c&a\end{smallmatrix}\right)\left(\begin{smallmatrix}a&b\\c&d\end{smallmatrix}\right)\).

\begin{example}\label{b2.19}
\(A=\left(\begin{smallmatrix}3&5\\1&2\end{smallmatrix}\right)\): \(ad-bc=6-5=1\), so
\(A^{-1}=\left(\begin{smallmatrix}2&-5\\-1&3\end{smallmatrix}\right)\) (Proposition~\ref{b2.18}).
\(B=\left(\begin{smallmatrix}2&6\\1&3\end{smallmatrix}\right)\): \(ad-bc=6-6=0\), not invertible; indeed
\(x=(3,-1)'\ne0\) has \(Bx=(6-6,\ 3-3)'=0\) (Theorem~\ref{b2.16}, (i)\(\Leftrightarrow\)(iii)).
\end{example}

\Needspace{15.5\baselineskip}
\begin{example}\label{b2.20}
Find the inverse of \(A=\left(\begin{smallmatrix}1&-1&1&-1\\-1&1&-2&1\\2&-1&2&-1\\1&-1&1&0\end{smallmatrix}\right)\),
and write \(A\) and \(A^{-1}\) as products of elementary matrices.
\end{example}
\textit{Solution.}
\begin{board}[Forward elimination of \([A\mid I_4]\):]
\[
\begin{gathered}
\left(\begin{array}{cccc|cccc}1&-1&1&-1&1&0&0&0\\-1&1&-2&1&0&1&0&0\\2&-1&2&-1&0&0&1&0\\1&-1&1&0&0&0&0&1\end{array}\right)
\xrightarrow[R_4\leftarrow R_4-R_1]{R_2\leftarrow R_2+R_1,\ R_3\leftarrow R_3-2R_1}
\left(\begin{array}{cccc|cccc}1&-1&1&-1&1&0&0&0\\0&0&-1&0&1&1&0&0\\0&1&0&1&-2&0&1&0\\0&0&0&1&-1&0&0&1\end{array}\right)
\\
\xrightarrow{R_2\leftrightarrow R_3}
\left(\begin{array}{cccc|cccc}1&-1&1&-1&1&0&0&0\\0&1&0&1&-2&0&1&0\\0&0&-1&0&1&1&0&0\\0&0&0&1&-1&0&0&1\end{array}\right)
\xrightarrow{R_3\leftarrow-R_3}
\left(\begin{array}{cccc|cccc}1&-1&1&-1&1&0&0&0\\0&1&0&1&-2&0&1&0\\0&0&1&0&-1&-1&0&0\\0&0&0&1&-1&0&0&1\end{array}\right)
\end{gathered}
\]
\end{board}
Entry \((2,2)\) is \(0\) after the first arrow: Step 2 interchanges rows \(2\), \(3\); Step 3 multiplies row \(3\) by
\(-1\).
\begin{board}[Backward elimination, clearing columns \(4\), \(3\), \(2\); Corollary~\ref{b2.17}(a):]
\[
\begin{gathered}
\xrightarrow[R_2\leftarrow R_2-R_4]{R_1\leftarrow R_1+R_4}
\left(\begin{array}{cccc|cccc}1&-1&1&0&0&0&0&1\\0&1&0&0&-1&0&1&-1\\0&0&1&0&-1&-1&0&0\\0&0&0&1&-1&0&0&1\end{array}\right)
\xrightarrow{R_1\leftarrow R_1-R_3}
\left(\begin{array}{cccc|cccc}1&-1&0&0&1&1&0&1\\0&1&0&0&-1&0&1&-1\\0&0&1&0&-1&-1&0&0\\0&0&0&1&-1&0&0&1\end{array}\right)
\\
\xrightarrow{R_1\leftarrow R_1+R_2}
\left(\begin{array}{cccc|cccc}1&0&0&0&0&1&1&0\\0&1&0&0&-1&0&1&-1\\0&0&1&0&-1&-1&0&0\\0&0&0&1&-1&0&0&1\end{array}\right),
\qquad
A^{-1}=\begin{pmatrix}0&1&1&0\\-1&0&1&-1\\-1&-1&0&0\\-1&0&0&1\end{pmatrix}
\end{gathered}
\]
\end{board}
Check: \(\row_3(A)A^{-1}=(2,-1,2,-1)A^{-1}=(0+1-2+1,\ 2+0-2+0,\ 2-1+0+0,\ 0+1+0-1)=(0,0,1,0)\); similarly for the
other rows.
\begin{board}[The nine operations; Theorems~\ref{b2.15}(c), \ref{b2.16} (last sentence), \ref{b2.13}(a),(b),
\ref{b2.15}(b):]
\boardtop
\begin{align*}
&E_1=T_{21}(1),\ E_2=T_{31}(-2),\ E_3=T_{41}(-1),\ E_4=P_{23},\ E_5=D_3(-1)\\
&E_6=T_{14}(1),\ E_7=T_{24}(-1),\ E_8=T_{13}(-1),\ E_9=T_{12}(1)\\
&E_9E_8\cdots E_1A=I_4\implies A^{-1}=E_9E_8\cdots E_1\\
&A=(A^{-1})^{-1}=E_1^{-1}E_2^{-1}\cdots E_9^{-1}
=T_{21}(-1)\,T_{31}(2)\,T_{41}(1)\,P_{23}\,D_3(-1)\,T_{14}(-1)\,T_{24}(1)\,T_{13}(1)\,T_{12}(-1)
\end{align*}
\end{board}

\begin{example}\label{b2.22}
The matrix \(A\) of Example~\ref{b2.10} is not invertible: an RREF of \(A\), the left part of the last matrix there (Definition~\ref{b2.2}), is
\(\left(\begin{smallmatrix}1&0&2&0\\0&1&-1&0\\0&0&0&1\\0&0&0&0\end{smallmatrix}\right)\ne I_4\). Accordingly \(Ax=0\)
has the nontrivial solution \(x=(-2,1,1,0)'\), and \(Ax=y\) has no solution for \(y=(1,0,0,0)'\)
(Theorem~\ref{b2.16}, (iii) and (iv)).
\end{example}

\subsection*{5. Solving \(Ax=b\)}

\begin{theorem}\label{b2.26}
Let \(A\in\R^{m\times n}\), \(b\in\R^m\), and let row operations take \([A\mid b]\) to \([H\mid d]\) in echelon form
(\(H\in\R^{m\times n}\)). Then \(Ax=b\) is consistent \(\iff\) \([H\mid d]\) has no row \((0\ \cdots\ 0\mid c)\) with
\(c\ne0\) \(\iff\) the last column of \([H\mid d]\) is not a pivot column. If row operations take \([A\mid b]\) to any
matrix with a row \((0\ \cdots\ 0\mid c)\), \(c\ne0\), then \(Ax=b\) has no solution. If \([H\mid d]\) has no such row,
let \(x_j\) be a \textbf{pivot unknown} if column \(j\) of \(H\) is a pivot column, a \textbf{free unknown} otherwise
(as in Proposition~\ref{b2.7}, now for \(H\)):
each choice of values of the free unknowns gives, by \textbf{back substitution} (solving the equations of the nonzero
rows of \(H\), from the last one up, for their pivot unknowns), exactly one solution, and every solution arises so.
\end{theorem}
\begin{proof}
\begin{itemize}
\item \(Ax=b\) and \(Hx=d\) have the same solutions (Theorem~\ref{b2.3}(b)). In an echelon form, a row
\((0\ \cdots\ 0\mid c)\), \(c\ne0\), is exactly a row whose pivot lies in the last column.
\item If \([B\mid z]\), reached from \([A\mid b]\), has such a row, the equation \(0=c\) of \(Bx=z\) has no solution;
nor has \(Ax=b\) (Theorem~\ref{b2.3}(b)). So one may stop at any stage of the elimination.
\item No such row in \([H\mid d]\): each zero row of \(H\) gives \(0=0\). Let rows \(1,\dots,r\) of \(H=(h_{ij})\) be
nonzero, with pivots in columns \(k_1<\dots<k_r\).
\item Fix the free unknowns. Row \(i\) of \(H\) is \(0\) before column \(k_i\) and \(1\) there, so for
\(i=r,r-1,\dots,1\) equation \(i\) reads \(x_{k_i}+\sum_{j>k_i}h_{ij}x_j=d_i\); each \(x_j\), \(j>k_i\), is free or
already found. So it determines \(x_{k_i}\), and \(x\) solves \(Hx=d\), hence \(Ax=b\).
\item The equations force the pivot unknowns: every solution arises from its free unknowns, exactly once.\qedhere
\end{itemize}
\end{proof}

\Needspace{13\baselineskip}
\begin{theorem}\label{b2.27}
Let \(A\in\R^{m\times n}\) and \(b\in\R^m\).
\begin{enumerate}
\item If \(Ax_p=b\) (a \textbf{particular solution}), the solutions of \(Ax=b\) are exactly the \(x_p+z\) with
\(Az=0\).
\item \(Ax=b\) has no solution, exactly one, or infinitely many. If it is consistent: exactly one \(\iff\) no free
unknowns (every column of \(H\) in Theorem~\ref{b2.26} is a pivot column).
\item If \(m=n\): \(Ax=b\) has exactly one solution \(\iff\) \(A\) is invertible; this solution is \(x=A^{-1}b\).
\end{enumerate}
\end{theorem}
\begin{proof}
\begin{board}[(a) \lec{1}, Theorem~3(a):]
\boardtop
\begin{align*}
Az=0&\implies A(x_p+z)=Ax_p+Az=b\\
Ax=b&\implies z=x-x_p:\ Az=b-b=0,\ x=x_p+z
\end{align*}
\end{board}
(b)
\begin{itemize}
\item Let \(Ax=b\) be consistent: forward elimination (Definition~\ref{b2.5}, Steps 1--4) takes \([A\mid b]\) to an
echelon form \([H\mid d]\), with no row \((0\ \cdots\ 0\mid c)\), \(c\ne0\) (Theorem~\ref{b2.26}).
\item By Theorem~\ref{b2.26}, each choice of the free unknowns gives exactly one solution, and every solution arises
so; different choices give different solutions (they differ in a free unknown, an entry of \(x\)).
\item No free unknowns: exactly one solution; a free unknown: each of its values gives another solution, infinitely
many.
\end{itemize}
\begin{board}[(c) \lec{1}, Theorem~3(a):]
\boardtop
\begin{align*}
A\text{ invertible}&\implies A(A^{-1}b)=b;\quad Ax=b\implies x=A^{-1}(Ax)=A^{-1}b\\
\text{one solution }x_p&\implies(Az=0\implies z=0)\quad\text{((a): else \(x_p+z\) is a second solution)}\\
&\implies A\text{ invertible}\quad\text{(Theorem~\ref{b2.16}, (iii)\(\Rightarrow\)(i))}
\end{align*}
\end{board}
\nopagebreak
\end{proof}

\Needspace{14.5\baselineskip}
\begin{example}\label{b2.28}
Solve \(x_1-x_2+x_3-x_4=1\), \(-x_1+x_2-2x_3+x_4=0\), \(2x_1-x_2+2x_3-x_4=3\), \(x_1-x_2+x_3=3\).
\end{example}
\textit{Solution.}
\begin{board}[\(A\) of Example~\ref{b2.20}, \(b=(1,0,3,3)'\); forward elimination of \([A\mid b]\) by the operations
there:]
\[
\begin{gathered}
\left(\begin{array}{cccc|c}1&-1&1&-1&1\\-1&1&-2&1&0\\2&-1&2&-1&3\\1&-1&1&0&3\end{array}\right)
\xrightarrow[R_4\leftarrow R_4-R_1]{R_2\leftarrow R_2+R_1,\ R_3\leftarrow R_3-2R_1}
\left(\begin{array}{cccc|c}1&-1&1&-1&1\\0&0&-1&0&1\\0&1&0&1&1\\0&0&0&1&2\end{array}\right)
\\
\xrightarrow{R_2\leftrightarrow R_3}
\left(\begin{array}{cccc|c}1&-1&1&-1&1\\0&1&0&1&1\\0&0&-1&0&1\\0&0&0&1&2\end{array}\right)
\xrightarrow{R_3\leftarrow-R_3}
\left(\begin{array}{cccc|c}1&-1&1&-1&1\\0&1&0&1&1\\0&0&1&0&-1\\0&0&0&1&2\end{array}\right)
\end{gathered}
\]
\end{board}
\begin{board}[Echelon form, no row \((0\ 0\ 0\ 0\mid c)\), \(c\ne0\): consistent (Theorem~\ref{b2.26}); no free
unknowns: one solution (Theorem~\ref{b2.27}(b)):]
\boardtop
\begin{align*}
&x_4=2,\quad x_3=-1,\quad x_2=1-x_4=-1,\quad x_1=1+x_2-x_3+x_4=1-1+1+2=3\\
&x=(3,-1,-1,2)';\quad\text{check with \(A^{-1}\) of Example~\ref{b2.20} (Theorem~\ref{b2.27}(c)):}\\
&A^{-1}b=(0+0+3+0,\ -1+0+3-3,\ -1+0+0+0,\ -1+0+0+3)'=(3,-1,-1,2)'
\end{align*}
\end{board}

\Needspace{11.5\baselineskip}
\begin{example}\label{b2.29}
Solve \(x_1+x_2+x_3-2x_4=2\), \(x_1+x_2+2x_3-3x_4=3\), \(2x_1+x_2+5x_3-4x_4=7\).
\end{example}
\textit{Solution.}
\begin{board}[\(A\) of Example~\ref{b2.6}, \(b=(2,3,7)'\); forward elimination of \([A\mid b]\) by the operations
there:]
\[
\begin{gathered}
\left(\begin{array}{cccc|c}1&1&1&-2&2\\1&1&2&-3&3\\2&1&5&-4&7\end{array}\right)
\xrightarrow[R_3\leftarrow R_3-2R_1]{R_2\leftarrow R_2-R_1}
\left(\begin{array}{cccc|c}1&1&1&-2&2\\0&0&1&-1&1\\0&-1&3&0&3\end{array}\right)
\\
\xrightarrow{R_2\leftrightarrow R_3}
\left(\begin{array}{cccc|c}1&1&1&-2&2\\0&-1&3&0&3\\0&0&1&-1&1\end{array}\right)
\xrightarrow{R_2\leftarrow-R_2}
\left(\begin{array}{cccc|c}1&1&1&-2&2\\0&1&-3&0&-3\\0&0&1&-1&1\end{array}\right)
\end{gathered}
\]
\end{board}
\begin{board}[Echelon form, no row \((0\ 0\ 0\ 0\mid c)\), \(c\ne0\): consistent (Theorem~\ref{b2.26}); pivot columns
\(1,2,3\), \(x_4=t\) free: infinitely many solutions (Theorem~\ref{b2.27}(b)):]
\boardtop
\begin{align*}
&x_3=1+t,\quad x_2=-3+3x_3=3t,\quad x_1=2-x_2-x_3+2x_4=2-3t-(1+t)+2t=1-2t\\
&x=(1,0,1,0)'+t(-2,3,1,1)'\quad(t\in\R)
\end{align*}
\end{board}
This is the particular solution \((1,0,1,0)'\) plus the solutions \(t(-2,3,1,1)'\) of \(Ax=0\) found in
Example~\ref{b2.6}, as Theorem~\ref{b2.27}(a) says.

\bigskip
\textbf{Practice.} The Exercises and Problems of the textbook, Chapter~2: row reduction and reduced row echelon forms,
reading off the solutions, inverses by Gauss--Jordan elimination, elementary matrices, and solving \(Ax=b\).

\end{document}
